\documentclass[letterpaper,10pt]{article}

\usepackage{latexsym}
\usepackage[empty]{fullpage}
\usepackage{titlesec}
\usepackage{enumitem}
\usepackage[hidelinks]{hyperref}
\usepackage{fancyhdr}
\usepackage[spanish]{babel}
\usepackage[utf8]{inputenc}
\usepackage[T1]{fontenc}
\usepackage{tabularx}
\input{glyphtounicode}

\usepackage[sfdefault]{roboto}

\pagestyle{fancy}
\fancyhf{}
\fancyfoot{}
\renewcommand{\headrulewidth}{0pt}
\renewcommand{\footrulewidth}{0pt}

\addtolength{\oddsidemargin}{-0.5in}
\addtolength{\evensidemargin}{-0.5in}
\addtolength{\textwidth}{1in}
\addtolength{\topmargin}{-.6in}
\addtolength{\textheight}{1.2in}

\urlstyle{same}
\raggedbottom
\raggedright
\setlength{\tabcolsep}{0in}
\pdfgentounicode=1

\titleformat{\section}{\large\bfseries\scshape\raggedright}{}{0em}{}[\titlerule]

\newcommand{\resumeItem}[1]{\item\small{#1}}
\newcommand{\resumeSubheading}[4]{
\vspace{-1pt}\item
  \begin{tabular*}{0.97\textwidth}[t]{l@{\extracolsep{\fill}}r}
    \textbf{#1} & #2 \\
    \textit{#3} & \textit{#4} \\
  \end{tabular*}\vspace{-7pt}
}
\newcommand{\resumeSubHeadingList}{\begin{itemize}[leftmargin=0.15in, label={}]}
\newcommand{\resumeSubHeadingListEnd}{\end{itemize}}
\newcommand{\resumeItemListStart}{\begin{itemize}[leftmargin=0.35in, label=\textbullet]}
\newcommand{\resumeItemListEnd}{\end{itemize}\vspace{-4pt}}

\begin{document}

\begin{center}
  \textbf{\Huge Luis Alfredo Cuamatzi Flores} \\ \vspace{2pt}
  \small Ingeniero en Inteligencia Artificial / Desarrollador de Software \\ \vspace{3pt}
  \href{mailto:mexboxluis1@gmail.com}{mexboxluis1@gmail.com} \;|\;
  México \;|\; Remoto / Disponible para reubicación \\ \vspace{2pt}
  \href{https://github.com/MexboxLuis}{github.com/MexboxLuis} \;|\;
  \href{https://www.linkedin.com/in/mexboxluis}{linkedin.com/in/mexboxluis} \;|\;
  \href{https://mexboxluis.github.io}{mexboxluis.github.io}
\end{center}

\section*{Resumen}
Ingeniero en Inteligencia Artificial y desarrollador de software con experiencia en aplicaciones móviles, servicios backend e integración de modelos de IA. He trabajado en proyectos reales combinando Android, APIs, automatización, procesamiento de datos, CRM, visión por computadora e interfaces orientadas a producto para conectar usuarios, datos y procesos de negocio.

\section{Experiencia}
\resumeSubHeadingList

\resumeSubheading
  {Desarrollador de Software e IA}{Jun 2025 -- Presente}
  {Net2Alliance}{Ciudad de México, México}
  \resumeItemListStart
    \resumeItem{Desarrollo full-stack e integración de IA para plataformas de telecomunicaciones y CRM, trabajando con FastAPI, Python, PHP, JavaScript, WebSockets y AWS SageMaker.}
    \resumeItem{Construcción de 2 módulos y 2 endpoints para análisis de sentimientos y generación automática de resúmenes en PDF a partir de transcripciones de llamadas mediante AWS SageMaker.}
    \resumeItem{Desarrollo de un cliente de mensajería sobre WebSockets para WhatsApp con notificaciones en tiempo real, soporte multimedia y transferencias entre agentes.}
    \resumeItem{Implementación de 4 módulos CRM con 4 endpoints, operaciones CRUD, control de acceso para 3 roles de usuario, auditoría y reportes administrativos.}
    \resumeItem{Participación en rediseños de UI/UX con Figma y seguimiento de tareas con Jira, GitHub y tableros Kanban.}
  \resumeItemListEnd

\resumeSubheading
  {Ingeniero Backend e IA -- Proyecto de Titulación / I+D}{Ene 2025 -- Dic 2025}
  {L4L -- Reconstrucción 3D de Prendas para Virtual Try-On}{Tlaxcala, México}
  \resumeItemListStart
    \resumeItem{Diseño de una arquitectura cliente--servidor para Virtual Try-On, integrando una aplicación Android en Jetpack Compose, un servidor Linux de inferencia y un servidor Windows para simulación física y conversión FBX--GLB.}
    \resumeItem{Desarrollo de un orquestador en FastAPI para coordinar captura de imagen, generación y confirmación de avatar, reconstrucción de prendas, simulación física y entrega de modelos 3D en formato GLB.}
    \resumeItem{Integración de un modelo de estimación de pose y forma corporal desde una sola imagen 2D (SMPLer-X), adaptando su salida para compatibilidad con el pipeline de simulación.}
    \resumeItem{Adaptación de un modelo Transformer de dos niveles para reconstrucción de patrones de costura desde una imagen 2D (Sewformer), evaluando variantes con poda, PTQ y QAT sobre métricas geométricas y de conectividad.}
    \resumeItem{Compresión del modelo final con una razón de compresión de 5.43x y reducción de latencia de 38.75\%, alcanzando 32.26 MB y 165.63 ms de inferencia en CPU.}
  \resumeItemListEnd

\resumeSubheading
  {Desarrollador Android -- Servicio Social}{Oct 2024 -- Ago 2025}
  {Instituto Politécnico Nacional -- Programa Institucional de Tutorías}{Tlaxcala, México}
  \resumeItemListStart
    \resumeItem{Registro de asistencia por QR con validación en tiempo real, detección de duplicados y control de sesiones activas.}
    \resumeItem{Sistema con 3 roles de usuario (Admin, Tutor Par y Asesor) y paneles diferenciados según permisos.}
    \resumeItem{Generador de horarios semanales con detección de traslapes entre aulas y tutores.}
    \resumeItem{Sincronización offline con Firestore y exportación a CSV para reportes institucionales.}
  \resumeItemListEnd

\resumeSubHeadingListEnd

\section{Educación}
\resumeSubHeadingList

\resumeSubheading
  {Ingeniería en Inteligencia Artificial}{Ago 2021 -- Ene 2026}
  {Instituto Politécnico Nacional}{Tlaxcala, México}

\resumeSubheading
  {Programa de Formación de Líderes}{Ago 2024}
  {Queen Mary University of London}{Londres, Reino Unido}

\resumeSubHeadingListEnd

\section{Proyectos}
\resumeSubHeadingList

\resumeSubheading
  {Lexfy}{Dic 2024}
  {Proyecto Android + IA}{\href{https://github.com/MexboxLuis/Lexfy}{GitHub}}
  \resumeItemListStart
    \resumeItem{Aplicación Android con OCR y generación de imágenes mediante una interfaz de chat.}
    \resumeItem{Cliente Android conectado por HTTP a un backend local en Python/Flask con EasyOCR, GOT-OCR2\_0 y Together AI.}
    \resumeItem{\textbf{Tecnologías:} Kotlin, Jetpack Compose, Python, Flask, Firebase, EasyOCR, GOT-OCR2\_0, Transformers.}
  \resumeItemListEnd

\resumeSubheading
  {SmartEMG Vision}{Nov 2024}
  {Prototipo Android para interacción asistida}{\href{https://github.com/MexboxLuis/SmartEMG-Vision}{GitHub}}
  \resumeItemListStart
    \resumeItem{Prototipo móvil que combina clasificación de señales EMG y detección de objetos para sugerir acciones según el contexto visual.}
    \resumeItem{Cliente Android con Jetpack Compose y CameraX conectado por HTTP a servidores locales en Python con YOLOv8 y TensorFlow/Keras.}
    \resumeItem{\textbf{Tecnologías:} Kotlin, Jetpack Compose, CameraX, Python, Flask, TensorFlow, Keras, YOLOv8.}
  \resumeItemListEnd

\resumeSubheading
  {SmartCrops}{Ago 2024}
  {Prototipo móvil para monitoreo de invernaderos}{\href{https://github.com/MexboxLuis/SmartCrops}{GitHub}}
  \resumeItemListStart
    \resumeItem{Prototipo para monitoreo de cultivos, simulación de variables de invernadero y asistencia mediante chatbot basado en Gemini.}
    \resumeItem{Aplicación de actividad única con Jetpack Compose, CameraX, OkHttp y Coil.}
    \resumeItem{\textbf{Tecnologías:} Kotlin, Jetpack Compose, Gemini API, CameraX, OkHttp, Coil.}
  \resumeItemListEnd

\resumeSubheading
  {CommuniSync}{May 2024}
  {Plataforma comunitaria Android}{\href{https://github.com/MexboxLuis/CommuniSync}{GitHub}}
  \resumeItemListStart
    \resumeItem{Aplicación Android nativa para participación comunitaria, incluyendo reportes ciudadanos, encuestas, foros, directorio local, notificaciones, perfiles de usuario y acceso de invitados.}
    \resumeItem{Integración de Firebase Authentication y Firestore para inicio de sesión con email, Google e invitado, almacenamiento de datos y sincronización en tiempo real.}
    \resumeItem{\textbf{Tecnologías:} Kotlin, Jetpack Compose, Material 3, Firebase Auth, Firestore, Coil.}
  \resumeItemListEnd

\resumeSubheading
  {EduStream API}{Jun 2023}
  {Plataforma web de gestión de contenidos educativos}{\href{https://github.com/MexboxLuis/EduStreamAPI}{GitHub}}
  \resumeItemListStart
    \resumeItem{Plataforma web para subir, categorizar y reproducir videos educativos por materia.}
    \resumeItem{Backend con arquitectura MVC, manejo de archivos multimedia y sincronización entre almacenamiento físico y base de datos.}
    \resumeItem{\textbf{Tecnologías:} Node.js, Express, MySQL, Handlebars, Bootstrap, Multer.}
  \resumeItemListEnd

\resumeSubHeadingListEnd

\section{Habilidades Técnicas}
\resumeSubHeadingList
  \resumeItem{\textbf{Lenguajes:} Kotlin, Python, JavaScript, TypeScript, PHP, SQL}
  \resumeItem{\textbf{Móvil:} Android nativo, Jetpack Compose, Android Views/XML, Material 3, CameraX, Room, WorkManager, OkHttp, Coil, ZXing}
  \resumeItem{\textbf{Backend:} FastAPI, Flask, Node.js, Express, WebSockets, REST APIs}
  \resumeItem{\textbf{IA / ML:} TensorFlow, Keras, PyTorch, Transformers, EasyOCR, YOLOv8, SMPLer-X, Gemini API, Together AI}
  \resumeItem{\textbf{Bases de datos / nube:} AWS SageMaker, Firestore, Firebase Realtime Database, Firebase Storage, MySQL, SQLite}
  \resumeItem{\textbf{Arquitectura y metodologías:} Agile / Scrum, Kanban, MVC, MVVM, arquitectura offline-first, diseño de APIs REST}
  \resumeItem{\textbf{Herramientas:} Git, GitHub, Jira, Figma, Postman, Android Studio, VS Code, Linux/CLI, SSH}
\resumeSubHeadingListEnd

\section{Idiomas}
\resumeSubHeadingList
  \resumeItem{\textbf{Español:} Nativo}
  \resumeItem{\textbf{Inglés:} B2 -- Competencia profesional}
\resumeSubHeadingListEnd

\end{document}
